% !PRINTED 650
\weberpage{650}
%<<CONT
wir ein von den $x, y, \ldots$ freies, mit $\varphi$ assoziiertes Funktional
$\psi = \varepsilon\varphi$ bestimmen. Da dann $\varphi/\psi$ ganz ist, so müssen nach 7.
die Koeffizienten von $\varphi$ alle durch $\psi$, und folglich auch alle
durch $\varphi$ teilbar sein. Was zu beweisen war.

Damit ist zugleich bewiesen:

9. \gesperrt{Ein holomorphes ganzes Funktional geht in
ein assoziiertes über, wenn die Variablen anders
bezeichnet werden.}

10. \gesperrt{Jedes Funktional $\bar\omega$ in $\bar\Omega$ geht durch Multiplikation
 mit einer Einheit in $\bar Z$ in ein holomorphes
Funktional über.}

Denn man kann zunächst $\bar\omega$ als Quotienten zweier ganzer
holomorpher Funktionale
\[
\omega = \frac{\psi}{\varphi}
% [sic] printed \omega = \frac{\psi}{\varphi} without the bar over \omega, the sentence represents \bar\omega and (6) continues with \bar\omega
\]
darstellen. Setzt man $N(\varphi) = \varphi\varphi'$, so ist, wie die Formel (4)
zeigt, in der $a_m E_m = \pm N(\varphi)$ ist, $\varphi'$ mit einem holomorphen
Funktional assoziiert, und es folgt:
\weq{6}{E_m \bar\omega = \frac{\psi \varphi'}{a_m}.}

11. \gesperrt{Zu einem ganzen Funktional $\bar\omega$ kann man ein
zweites $\mu$ so wählen, daß $\mu$ durch beliebig gegebene
Primfunktionale $\pi_1, \pi_2, \ldots$ nicht teilbar wird, und daß}
\weq{7}{\bar\omega \mu = \alpha}
\gesperrt{eine Funktion in $\Omega$ wird.}

Man stelle $\bar\omega$ nach 10. durch ein holomorphes Funktional $\varphi$
dar. Dessen Koeffizienten sind nach 8. alle durch $\bar\omega$, aber nicht
alle durch $\bar\omega \pi_1$ teilbar. Es gibt also eine Funktion $\alpha_1$, die durch
$\bar\omega$, aber nicht durch $\bar\omega \pi_1$ teilbar ist.

Nun setze man:
\[
\bar\omega_1 = \bar\omega \pi_2 \pi_3 \ldots, \quad \bar\omega_2 = \bar\omega \pi_1 \pi_3 \ldots, \quad \bar\omega_3 = \omega \pi_2 \pi_3 \ldots
% [sic] printed \bar\omega_3 = \omega \pi_2 \pi_3 \ldots, the bar over \omega is missing and the factors repeat those of \bar\omega_1; the construction requires \bar\omega_3 = \bar\omega \pi_1 \pi_2 \pi_4 \ldots
\]
und bestimme:
\[
\begin{aligned}
&\alpha_1 \text{ teilbar durch } \bar\omega_1, \text{ nicht teilbar durch } \bar\omega_1 \pi_1, \\
&\alpha_2 \text{ \textquotedbl} \text{ \textquotedbl} \bar\omega_2, \text{ \textquotedbl} \text{ \textquotedbl} \text{ \textquotedbl} \bar\omega_2 \pi_2, \\
&\alpha_3 \text{ \textquotedbl} \text{ \textquotedbl} \bar\omega_3, \text{ \textquotedbl} \text{ \textquotedbl} \text{ \textquotedbl} \bar\omega_3 \pi_3.
\end{aligned}
\]
\[
\cdots\cdots\cdots\cdots\cdots\cdots\cdots\cdots\cdots\cdots\cdots\cdots\cdots\cdots\cdots\cdots\cdots
\]
Dann genügt
\[
\alpha = \alpha_1 + \alpha_2 + \alpha_3 + \cdots
\]
den Forderungen des Satzes 11.
